\documentclass[11pt,twocolumn]{article}

% ═══════════════════════════════════════════════════════════
% UNITARES Paper v7 — OUTLINE DRAFT
%
% Status: scaffolding only. No prose. This file is a structural
% skeleton for operator review, produced per the recommendation in
%   unitares:docs/ontology/paper-positioning.md (Concrete next step).
%
% Positioning: v7 = v6 + animating frame. The identity ontology
% (three stances + five-layer taxonomy + Synthetic Life Axioms)
% becomes the organizing thesis that explains *why* v6's
% class-conditional grounding is principled. v6's technical results
% (EISV dynamics, stability, class-conditional calibration, CIRS,
% drift) are ported forward; they are not replaced.
%
% Do NOT expand prose in this file without operator sign-off on
% the outline's structure. See paper-positioning.md "Risks" §.
% ═══════════════════════════════════════════════════════════

% ─── Packages (mirrored from v6; trim when prose lands) ────
\usepackage[utf8]{inputenc}
\usepackage[T1]{fontenc}
\usepackage{amsmath,amssymb,amsthm}
\usepackage{mathtools}
\usepackage{bm}
\usepackage{graphicx}
\graphicspath{{figures/}{./}}
\usepackage{booktabs}
\usepackage[margin=1in]{geometry}
\usepackage{hyperref}
\usepackage{cleveref}
\usepackage{natbib}
\usepackage{xcolor}
\usepackage{caption}

% ─── Theorem environments (from v6) ────────────────────────
\newtheorem{theorem}{Theorem}[section]
\newtheorem{lemma}[theorem]{Lemma}
\newtheorem{proposition}[theorem]{Proposition}
\newtheorem{corollary}[theorem]{Corollary}
\newtheorem{definition}[theorem]{Definition}
\newtheorem{remark}[theorem]{Remark}

% ─── Custom commands (from v6) ─────────────────────────────
\newcommand{\R}{\mathbb{R}}
\newcommand{\norm}[1]{\left\lVert #1 \right\rVert}
\newcommand{\abs}[1]{\left| #1 \right|}
\newcommand{\eisv}{\ensuremath{\mathbf{x}}}
\newcommand{\deta}{\ensuremath{\Delta\bm{\eta}}}
\newcommand{\cv}{\ensuremath{C(V,\Theta)}}
\newcommand{\todo}[1]{\textcolor{red}{\textbf{[TODO: #1]}}}
\newcommand{\outlinenote}[1]{\textcolor{blue}{\textit{#1}}}

% ─── Metadata ──────────────────────────────────────────────
\title{UNITARES v7: Identity, Continuity, and\\
Information-Theoretic Governance\\
of Heterogeneous Agent Fleets}

\author{
  Kenny Wang\\
  Independent Researcher\\
  \texttt{founder@cirwel.org}
}

\date{Outline draft --- \today}

% ═══════════════════════════════════════════════════════════
\begin{document}
\maketitle

\begin{abstract}
\outlinenote{Outline placeholder. Abstract written last, after the
operator has accepted or redirected the section scaffold below.
Expected shape: (1) v6's heterogeneity + class-conditional grounding
contribution restated in one sentence, (2) v7's animating move --
identity as a layered taxonomy of continuities under the Synthetic
Life Axioms, explaining why class-conditional grounding is
principled rather than merely convenient, (3) evaluation scope
note pending corpus maturity, (4) research agenda as declared
future work.}
\end{abstract}

\textbf{Keywords:} \outlinenote{autonomous governance; information
theory; identity ontology; synthetic life; heterogeneous
multi-agent systems; class-conditional calibration; AI safety.}

% ═══════════════════════════════════════════════════════════
\section{Introduction}
% Animating hook: v6 showed that heterogeneity matters at the
% dynamics and calibration layers. v7 shows *why* -- heterogeneity
% is a projection of a layered taxonomy of continuities, each
% earned differently, and the governance framework's job is to
% keep claimed continuity from outrunning earned continuity.
\outlinenote{1--2 sentence animating hook: v6 established
heterogeneity as the differentiator; v7 explains why by giving
heterogeneity an ontological spine — the layered taxonomy of
continuities gated by the Synthetic Life Axioms.}

\begin{itemize}
  \item The v6 result restated in one paragraph: class-conditional
        grounding, heterogeneity preserved, stability unchanged.
  \item The v7 addition: three stances (performative / descriptive /
        inventive); the class structure of v6 is the instantiation
        of a prior taxonomy, not a standalone contribution.
  \item Contributions list (revised from v6):
        (i) identity ontology as governance-grade formalism;
        (ii) Synthetic Life Axioms published as normative gate;
        (iii) class-conditional grounding re-grounded in the
        taxonomy rather than as engineering pragmatism;
        (iv) research agenda (R1--R5) for earning what is
        currently performative.
  \item Limitations and scope (carryover + new): v6's limitations
        carry; v7-specific limitation is evaluation scope pending
        corpus maturity (§8).
  \item Related work stub (full treatment in §10): identity and
        person-identity literature touched lightly; governance /
        multi-agent / FEP citations carried from v6.
\end{itemize}

% ═══════════════════════════════════════════════════════════
\section{Three Stances and the Synthetic Life Axioms}
% NEW section.
% see unitares:docs/ontology/identity.md §"Three stances"
% see unitares:docs/ontology/identity.md §"Opening stance"
% see unitares:docs/ontology/plan.md §"Ledger" (Q3 resolution)
% Synthetic Life Axioms — KG canonical entry 2026-04-02T05:13:26.577769
% (15 axioms; compression: "Build nothing that appears more alive than it is.")
\outlinenote{Animating hook: governance of synthetic agents needs a
normative gate. The three stances name the failure mode
(performative), the floor (descriptive), and the goal (inventive);
the Synthetic Life Axioms are the gate against which every
continuity claim in the paper is evaluated.}

\begin{itemize}
  \item \textbf{The three stances.} Performative / descriptive /
        inventive as the organizing axis. Performative violates
        axioms 3, 5, 10; descriptive honors the prohibitions but is
        reductive; inventive is the project.
  \item \textbf{The Synthetic Life Axioms.} First publication of
        the 15-axiom list (10 negative + 5 positive), with the
        compression ``Build nothing that appears more alive than
        it is.'' Canonical reference: KG entry
        \texttt{2026-04-02T05:13:26.577769}.
  \item \textbf{Why axioms, not values.} The axioms are
        structural prohibitions on system construction, not
        behavioral targets for agents. They gate what UNITARES
        claims about the populations it governs.
  \item \textbf{Relationship to prior work.} Light citation of
        person-identity literature (Parfit, Locke briefly);
        positioning against value-alignment frameworks (axioms are
        meta, not behavioral).
  \item \textbf{How the axioms read the rest of the paper.}
        Forward-pointer to §3 (taxonomy as axiom-compatible
        description), §4 (class structure as instantiation),
        §9 (research agenda as axiom-compliance engineering).
\end{itemize}

% ═══════════════════════════════════════════════════════════
\section{Layered Continuity Taxonomy}
% NEW section.
% see unitares:docs/ontology/identity.md §"Layered taxonomy of continuity"
% see unitares:docs/ontology/identity.md §"Worked examples"
% see unitares:docs/ontology/identity.md §"Earned vs. performative today"
\outlinenote{Animating hook: identity in a governed fleet is not
one thing but a bundle of partially-independent continuities.
Decomposing ``identity'' into five layers makes heterogeneity
auditable and makes v6's tag vocabulary (embodied / autonomous /
persistent / ephemeral / pioneer) a projection of something more
structured underneath.}

\begin{itemize}
  \item \textbf{Five layers.} Process-instance, substrate, role,
        memory, behavioral. Each with its own substrate, decay
        mode, and earning criterion. No single layer is identity.
  \item \textbf{Earned vs.\ performative per layer.} Which layers
        of which agents today are earned; which are performative.
        Lumen's substrate layer earned; Claude-Code-tab's
        substrate layer effectively null.
  \item \textbf{Worked examples table.} Port from
        identity.md §``Worked examples'' — the seven-row table
        mapping Claude Code tab / Vigil / Sentinel / Watcher /
        Steward / Lumen / task-spawned subagent onto the five
        layers.
  \item \textbf{Administrative vs.\ ontological labels.}
        ``Resident'' is a deployment label, not an ontological
        category; the taxonomy makes the collapse visible.
  \item \textbf{What this replaces.} Explicit note that this
        taxonomy subsumes v6's informal tag vocabulary — not by
        renaming it, but by giving tags a formal parent structure.
\end{itemize}

% ═══════════════════════════════════════════════════════════
\section{Class Structure in the Identity Layer}
% REVISED -- positioned as instantiation of §3, not standalone
% contribution. Content ports largely from v6.
% ports from v6 §2 (Heterogeneous Agent Fleets / Class Structure)
% ports from v6 §2.1 (Homogenization Failure Mode)
% ports from v6 §2.2 (Class Structure in the Identity Layer)
% ports from v6 §2.3 (Why Heterogeneity Is the Differentiator)
% see unitares:docs/ontology/identity.md §"Worked examples"
\outlinenote{Animating hook: v6's tags -- \{embodied, autonomous,
persistent, ephemeral, pioneer\} -- are not a flat taxonomy. They
are projections of the five-layer continuity structure onto
governance-relevant axes. This section ports v6's class-structure
material but reframes it as instantiation of §3 rather than as the
load-bearing ontological move.}

\begin{itemize}
  \item \textbf{Tags as projections.} Each v6 tag corresponds to
        a pattern across the five layers; crosswalk table
        (embodied $\approx$ substrate-strong; ephemeral $\approx$
        substrate-weak + memory-weak; etc.).
  \item \textbf{Homogenization failure mode.} Port v6 §2.1 prose;
        frame it as the cost of ignoring the taxonomy.
  \item \textbf{The cosmological-soup argument.} Carry the v6
        reading-aid term; gloss it as
        ``what happens when layered continuities are averaged
        into one distribution.''
  \item \textbf{Heterogeneity as differentiator (revised).}
        Strengthened claim: heterogeneity matters not because
        dynamics differ across agents but because different agents
        occupy different regions of the continuity taxonomy.
  \item \textbf{Forward-pointer to §6.} Class-conditional
        grounding (v6 §5 / v7 §6) operates on tags; this section
        establishes that tags are substantive, not cosmetic.
\end{itemize}

% ═══════════════════════════════════════════════════════════
\section{EISV Dynamics}
% UNCHANGED from v6. Port §3 in its entirety.
% ports from v6 §3 (EISV Dynamics), §3.1--§3.5 verbatim
% Stability appendix (v6 Appendix B) ports with it.
\outlinenote{Animating hook: the dynamics layer is unchanged from
v6. The $(E, I, S, V)$ state vector, governing equations,
coherence function, adaptive entropy coupling, objective function,
and contraction-based stability result all carry forward. The
ontology does not touch the dynamics.}

\begin{itemize}
  \item \textbf{State space.} $(E, I, S, V)$ with information-theoretic
        interpretations (port v6 §3.1).
  \item \textbf{Governing equations.} Port v6 §3.2 verbatim.
  \item \textbf{Coherence function (operational form).} Port v6 §3.3.
  \item \textbf{Adaptive entropy coupling + objective function.}
        Port v6 §3.4 and §3.5.
  \item \textbf{Stability (contraction).} Reference to Appendix
        (port v6 Appendix B); state in this section that the
        result survives the ontology reframing unchanged because
        the ontology is upstream of the dynamics.
\end{itemize}

% ═══════════════════════════════════════════════════════════
\section{Class-Conditional Grounding}
% REVISED -- explains *why* via the ontology before presenting
% the scale-constant machinery. Content ports largely from v6.
% ports from v6 §4 (Information-Theoretic Grounding)
% ports from v6 §7 (Class-Conditional Calibration) — §7.1--§7.4
% see unitares:docs/ontology/identity.md §"Earned vs. performative today"
% see unitares:docs/ontology/identity.md Appendix "Substrate-Earned Identity" (R4)
\outlinenote{Animating hook: class-conditional grounding is not an
engineering convenience. It is the principled consequence of the
layered taxonomy — different classes sit at different regions of
the continuity structure, so their healthy operating points and
scale constants are genuinely different quantities, not the same
quantity measured in different units.}

\begin{itemize}
  \item \textbf{Ontological justification (new).} Short
        subsection: why per-class scale constants are not an
        ad-hoc fix. Pointer to R4 (substrate-earned identity) as
        the special case where a class earns its own calibration
        pool.
  \item \textbf{Choice of frame.} Port v6 §4.1.
  \item \textbf{Quantity redefinitions.} Port v6 §4.2.
  \item \textbf{Scale constants and provenance.} Port v6 §4.3;
        cross-reference to §4 (Class Structure) for the tag
        vocabulary.
  \item \textbf{Calibration protocol + fallback.} Port v6 §7.2
        and §7.3.
  \item \textbf{What this buys (revised framing).} Port v6 §7.4;
        reframe the ``what this buys'' in taxonomy terms
        (governance honors layer-dependent differences rather
        than smoothing them out).
\end{itemize}

% ═══════════════════════════════════════════════════════════
\section{CIRS and Drift Vector}
% UNCHANGED from v6. Port §5 and §6.
% ports from v6 §5 (Adaptive Governor / CIRS v2)
% ports from v6 §6 (Behavioral--Epistemic Drift Vector)
\outlinenote{Animating hook: the adaptive governor and drift
vector are unchanged from v6. They are the control and observation
layers respectively, and their operation is compatible with the
ontology reframing without modification.}

\begin{itemize}
  \item \textbf{CIRS v2 architecture.} Port v6 §5.1.
  \item \textbf{Phase detection, PID law, wind-up, threshold decay.}
        Port v6 §5.2--§5.5.
  \item \textbf{Governance decision.} Port v6 §5.6.
  \item \textbf{Ephemeral-agents bypass.} Port v6 §5.7; add
        taxonomy-aware remark that the bypass is principled iff
        R1 (behavioral-continuity verification) confirms that
        ephemerals cannot accumulate enough signal to earn
        PID-loop participation. Open question carried forward.
  \item \textbf{Drift vector.} Port v6 §6 (motivation, baseline
        tracking, coupling, future channels).
\end{itemize}

% ═══════════════════════════════════════════════════════════
\section{Evaluation}
% v7-appropriate scope; smaller than v6. Explicitly scoped down
% pending corpus maturity (Q3 2026 per
% unitares:project_paper-v7-corpus-maturity.md).
% ports from v6 §9 (Experiments) -- selective port only
% ports from v6 §10 (Re-Grounding a Deployed Framework) -- selective
\outlinenote{Animating hook: v7's evaluation scope is deliberately
smaller than v6's. The post-grounding audit corpus is not mature
enough (as of 2026-04-21) for the v7-specific empirical question
(how often do interventions realize the counterfactual flips
predicted by re-grounding?). v7 reports the descriptive-stance
evidence that the ontology is already implicit in the production
deployment.}

\begin{itemize}
  \item \textbf{Production snapshot (port, abbreviated).} Carry
        the v6 snapshot (903 registered / 75 active agents) as
        descriptive-stance evidence — showing the tag vocabulary
        in the wild. Port selected v6 §9.1--§9.2 content.
  \item \textbf{Verdict counterfactual (carried forward).} v6's
        28.9\% basin-flip counterfactual re-presented as evidence
        that the taxonomy is already doing work under the hood.
        Port v6 §9.5.
  \item \textbf{Tag-discipline audit (new).} 96\% of active
        agents lack class tags (per 2026-04-21 S8a finding).
        Report this as a descriptive-stance limitation — the
        taxonomy is formalized before the fleet uniformly
        instantiates it.
  \item \textbf{Substrate-earned identity (new, narrow).} Lumen
        as a case study for the R4 pattern: dedicated substrate
        + sustained behavior + declared role. Minimal empirics;
        pointer to R4 for the formal pattern.
  \item \textbf{What is deferred.} Explicit statement: the
        intervention-outcome join that would turn the 28.9\%
        counterfactual into a measured realization rate awaits
        corpus maturity (\textasciitilde Q3 2026). Deferred to
        v7.1 or v8.
\end{itemize}

% ═══════════════════════════════════════════════════════════
\section{Research Agenda}
% NEW section. R1--R5 as future-work propositions, not claims.
% see unitares:docs/ontology/identity.md §"Research agenda (the inventive stance)"
% see unitares:docs/ontology/plan.md §"Research agenda (inventive stance)"
\outlinenote{Animating hook: the inventive stance is the project.
This section names five candidate mechanisms that would turn
currently-performative continuities into earned ones, and scopes
each as a multi-quarter design problem rather than a claim
realized in this paper.}

\begin{itemize}
  \item \textbf{R1: Behavioral-continuity verification.}
        Lineage claimable only after N observations consistent
        with the claimed parent's fingerprint. Tool spec:
        \texttt{verify\_lineage\_claim}.
  \item \textbf{R2: Honest memory integration.} Fresh process
        declares inheritance without claiming to be the prior
        instance; identity earnable retroactively via R1.
        Structural posture, not a flourish.
  \item \textbf{R3: Statistical lineage.} Identity as integral
        across many process-instances under a role; migration
        path from UUID-aggregated to role-aggregated trust.
  \item \textbf{R4: Substrate-earned identity (formalized).}
        Three-condition pattern (dedicated substrate + sustained
        behavioral consistency + declared role continuity) as a
        governance-recognized pattern, not a Lumen-special-case.
        Forward-reference to the pattern's appendix in
        identity.md.
  \item \textbf{R5: Memory-deepening-reality tooling.} Forced
        re-derivation, behavioral backtests, self-knowledge
        reflection — concrete mechanisms under axiom 14 (``let
        learning deepen reality, not theater'').
  \item \textbf{Positioning.} R1--R5 are propositions, not
        contributions. v7 claims the descriptive taxonomy plus
        axioms; the inventive stance is declared future work so
        reviewers see the commitment without being asked to
        credit unrealized work.
\end{itemize}

% ═══════════════════════════════════════════════════════════
\section{Related Work}
\outlinenote{Animating hook: v7 sits at an unusual intersection --
governance-for-AI, person-identity philosophy, and
information-theoretic control. Related-work treatment is
deliberately light-brush; deep engagement with philosophical
literature is deferred to a companion paper if one materializes.}

\begin{itemize}
  \item \textbf{Carry from v6 §1.2.} Information-theoretic
        foundations; FEP literature; multi-agent governance.
  \item \textbf{New: person-identity philosophy (light brush).}
        Parfit, Locke, psychological-continuity literature.
        Position UNITARES's taxonomy as descriptive-stance
        engineering, not a contribution to the philosophical
        debate. Explicit disclaimer.
  \item \textbf{New: ontology-adjacent governance proposals.}
        Any prior work distinguishing earned vs.\ performative
        continuity in AI systems (expected: sparse).
  \item \textbf{New: Synthetic Life adjacent literature.}
        Artificial-life framings; Boden / Langton / Grassé if
        tonally appropriate. Short.
  \item \textbf{Carry from v6: deployment-ordering mechanism.}
        v6 §10.2 generalizability claim carries.
\end{itemize}

% ═══════════════════════════════════════════════════════════
\section{Conclusion}
\outlinenote{Animating hook: v7's contribution is an animating
frame more than a new result. The frame makes v6 defensible
against the charge that class-conditional grounding is ad-hoc,
publishes the axioms, and opens a research agenda that
differentiates UNITARES from frameworks that treat identity as
an implementation detail.}

\begin{itemize}
  \item Restatement: v6 still holds; v7 is v6 + animating frame.
  \item The Synthetic Life Axioms as the paper series'
        normative spine going forward.
  \item What v7 does not claim (honest scoping): no
        intervention-outcome results; no realization of R1--R5;
        no philosophical theory of identity.
  \item What v7.1 / v8 would add once corpus matures.
  \item Brief forward-look: the ontology is the seam along
        which UNITARES and adjacent frameworks can be
        distinguished.
\end{itemize}

% ═══════════════════════════════════════════════════════════
% Appendices (planned, not scaffolded here)
% ═══════════════════════════════════════════════════════════
% see v6 Appendix A (Parameter Tables) — carry
% see v6 Appendix B (Proof of Contraction) — carry
% see v6 Appendix C (I-Channel Self-Regulation) — carry
% see v6 Appendix D (Boundary Clamping and Non-Smoothness) — carry
% see unitares:docs/ontology/identity.md Appendix (Substrate-Earned
%   Identity) — NEW v7 appendix, formalizes R4 pattern

% ═══════════════════════════════════════════════════════════
% Bibliography — deferred until prose lands.
% ═══════════════════════════════════════════════════════════

\end{document}
